% ============================================================
% FUTA FRESHERS' SURVIVAL GUIDE — LaTeX Source
% Author: Abideen Yakub Opeyemi (Ó'YẸMÍ)
% Version: 2.2.0
% Compiles to: futa-guide.pdf
% ============================================================

\documentclass[11pt,a4paper]{article}

% ============================================================
% PACKAGES
% ============================================================
\usepackage[utf8]{inputenc}
\usepackage[T1]{fontenc}
\usepackage[margin=2.5cm, top=2cm, bottom=2.5cm]{geometry}
\usepackage{xcolor}
\usepackage{titlesec}
\usepackage{enumitem}
\usepackage{tcolorbox}
\usepackage{tikz}
\usepackage{hyperref}
\usepackage{fancyhdr}
\usepackage{fontawesome5}
\usepackage{graphicx}
\usepackage{parskip}
\usepackage{multicol}
\usepackage{array}
\usepackage{longtable}
\usepackage{booktabs}
\usepackage{tocloft}
\usepackage{sectsty}
\usepackage{microtype}
\usepackage{lmodern}

\tcbuselibrary{skins,breakable}

% ============================================================
% COLORS (matching the website theme)
% ============================================================
\definecolor{primary}{HTML}{003366}
\definecolor{secondary}{HTML}{990000}
\definecolor{accent}{HTML}{FFCC00}
\definecolor{teal}{HTML}{00695C}
\definecolor{gpaDark}{HTML}{001A33}
\definecolor{successGreen}{HTML}{28A745}
\definecolor{warningYellow}{HTML}{FFC107}
\definecolor{dangerRed}{HTML}{DC3545}
\definecolor{lightGray}{HTML}{F8F9FA}
\definecolor{midGray}{HTML}{E9ECEF}
\definecolor{textGray}{HTML}{495057}
\definecolor{inspirational}{HTML}{E3F2FD}
\definecolor{inspirationalTeal}{HTML}{E0F2F1}
\definecolor{warningBg}{HTML}{FEF3C7}
\definecolor{warningBorder}{HTML}{F59E0B}
\definecolor{warningText}{HTML}{B45309}
\definecolor{glossaryBg}{HTML}{FEF9E7}

% ============================================================
% HYPERREF SETUP
% ============================================================
\hypersetup{
    colorlinks=true,
    linkcolor=primary,
    urlcolor=primary,
    citecolor=primary,
    pdftitle={FUTA Freshers' Survival Guide},
    pdfauthor={Abideen Yakub Opeyemi},
    pdfsubject={Complete First Semester Guide},
    pdfkeywords={FUTA, freshers, survival guide, 100 level}
}

% ============================================================
% HEADER & FOOTER
% ============================================================
\pagestyle{fancy}
\fancyhf{}
\fancyhead[L]{\textcolor{primary}{\textbf{FUTA Freshers' Survival Guide}}}
\fancyhead[R]{\textcolor{textGray}{\small 2026/2027 Session}}
\fancyfoot[C]{\thepage}
\fancyfoot[L]{\textcolor{textGray}{\small By O'yemi}}
\fancyfoot[R]{\textcolor{textGray}{\small Not officially affiliated with FUTA}}
\renewcommand{\headrulewidth}{0.5pt}
\renewcommand{\footrulewidth}{0.3pt}
\renewcommand{\headrule}{\hbox to\headwidth{\color{primary}\leaders\hrule height \headrulewidth\hfill}}

% ============================================================
% SECTION STYLING
% ============================================================
\titleformat{\section}
    {\color{primary}\Huge\bfseries}
    {\thesection}{1em}{}
    [\vspace{0.3em}\color{primary}\hrule\vspace{0.5em}]

\titleformat{\subsection}
    {\color{primary}\Large\bfseries}
    {\thesubsection}{1em}{}

\titleformat{\subsubsection}
    {\color{primary}\large\bfseries}
    {\thesubsubsection}{1em}{}

\titlespacing*{\section}{0pt}{2em}{1em}
\titlespacing*{\subsection}{0pt}{1.5em}{0.8em}

% ============================================================
% CUSTOM BOXES
% ============================================================
\newtcolorbox{inspirational}[1][]{
    colback=inspirational, colframe=primary,
    leftrule=4pt, rightrule=0pt, toprule=0pt, bottomrule=0pt,
    arc=4pt, boxsep=8pt, left=10pt, right=10pt, top=8pt, bottom=8pt,
    breakable, #1
}

\newtcolorbox{warningbox}[1][]{
    colback=warningBg, colframe=warningBorder,
    leftrule=4pt, rightrule=0pt, toprule=0pt, bottomrule=0pt,
    arc=4pt, boxsep=8pt, left=10pt, right=10pt, top=8pt, bottom=8pt,
    breakable, #1
}

\newtcolorbox{successbox}[1][]{
    colback=successGreen!15, colframe=successGreen,
    leftrule=4pt, rightrule=0pt, toprule=0pt, bottomrule=0pt,
    arc=4pt, boxsep=8pt, left=10pt, right=10pt, top=8pt, bottom=8pt,
    breakable, #1
}

\newtcolorbox{moneybox}[1][]{
    colback=warningYellow!25, colframe=warningYellow,
    boxrule=0.5pt, arc=4pt,
    boxsep=8pt, left=10pt, right=10pt, top=8pt, bottom=8pt,
    breakable, #1
}

\newtcolorbox{gpabox}[1][]{
    colback=gpaDark, colframe=accent, coltext=white,
    boxrule=0pt, leftrule=6pt, arc=4pt,
    boxsep=10pt, left=15pt, right=15pt, top=12pt, bottom=12pt,
    breakable, #1
}

\newtcolorbox{quotebox}[1][]{
    colback=lightGray, colframe=primary,
    boxrule=0pt, leftrule=3pt, arc=4pt,
    boxsep=6pt, left=12pt, right=12pt, top=6pt, bottom=6pt,
    breakable, #1
}

\newtcolorbox{mindsetbox}[1][]{
    colback=white, colframe=primary,
    boxrule=0.5pt, leftrule=5pt, arc=4pt,
    boxsep=10pt, left=15pt, right=15pt, top=10pt, bottom=10pt,
    breakable, #1
}

\newtcolorbox{coursebox}[1][]{
    colback=lightGray, colframe=primary,
    boxrule=0.5pt, leftrule=4pt, arc=4pt,
    boxsep=8pt, left=12pt, right=12pt, top=10pt, bottom=10pt,
    breakable, #1
}

% ============================================================
% CUSTOM COMMANDS
% ============================================================
\newcommand{\quotetext}[1]{\begin{quotebox}\itshape\color{primary}#1\end{quotebox}}
\newcommand{\highlight}[1]{\textbf{\textcolor{primary}{#1}}}
\newcommand{\keyterm}[1]{\textbf{\textcolor{secondary}{#1}}}

\renewcommand{\cftsecfont}{\color{primary}\bfseries}
\renewcommand{\cftsubsecfont}{\color{textGray}}

% ============================================================
% DOCUMENT START
% ============================================================
\begin{document}

% ============================================================
% COVER PAGE
% ============================================================
\begin{titlepage}
    \centering

    \begin{tikzpicture}[remember picture, overlay]
        \fill[primary] (current page.north west) rectangle ([yshift=-6cm]current page.north east);
        \fill[accent] ([yshift=-6cm]current page.north west) rectangle ([yshift=-6.4cm]current page.north east);
    \end{tikzpicture}

    \vspace*{1cm}

    {\color{white}\Huge\bfseries FEDERAL UNIVERSITY\\[0.3em] OF TECHNOLOGY, AKURE\par}

    \vspace{1cm}

    {\color{accent}\rule{10cm}{2pt}\par}

    \vspace{0.5cm}

    {\color{white}\Large\itshape Freshers' Survival Guide\par}

    \vspace{2cm}

    {\color{primary}\Huge\bfseries FUTA\\[0.3em] FRESHERS'\\[0.3em] SURVIVAL GUIDE\par}

    \vspace{0.5cm}

    {\color{primary}\Large Complete First Semester Guide\par}

    \vspace{0.3cm}

    {\color{textGray}\large 2026/2027 Academic Session\par}

    \vspace{2cm}

    \begin{tcolorbox}[colback=lightGray, colframe=primary, width=12cm, arc=6pt, boxsep=10pt]
        \centering
        {\color{primary}\textbf{Compiled by:}}\\[0.3em]
        {\Large\bfseries Abideen Yakub Opeyemi}\\[0.2em]
        {\itshape (Ó'YẸMÍ)}\\[0.5em]
        {\small Department of Electrical and Electronics Engineering}\\[0.3em]
        {\small\faPhone\ +234 705 037 6975}\\[0.2em]
        {\small\faEnvelope\ abideenyakub08@gmail.com}
    \end{tcolorbox}

    \vfill

    \begin{quotebox}
        \centering\itshape\color{primary}
        "The future belongs to those who believe in the beauty of their dreams."\\
        \textnormal{\small --- Eleanor Roosevelt}
    \end{quotebox}

    \vspace{1cm}

    {\color{textGray}\footnotesize This is an independent student-created resource. It is \textbf{not} an official FUTA publication.\par}

    \vspace{0.5cm}

    {\color{textGray}\footnotesize \copyright\ 2026 FUTA Survival Guide. All rights reserved.\par}

\end{titlepage}

% ============================================================
% TABLE OF CONTENTS
% ============================================================
\newpage
\tableofcontents

% ============================================================
% SECTION: WELCOME MESSAGE (NEW PAGE)
% ============================================================
\newpage
\section*{A Welcome Message}
\addcontentsline{toc}{section}{A Welcome Message}

To the brilliant minds who have earned their place here, \textbf{welcome to the Federal University of Technology, Akure (FUTA)!}

This achievement is a testament to your dedication, and you should be immensely proud. You are now part of a legacy of excellence and innovation. \highlight{You are a FUTArian.}

\begin{inspirational}
    \textbf{\color{primary}\Large 🏆 You Are Architects of Your Future}

    \vspace{0.5em}

    The journey you're beginning will be one of the most defining chapters of your life. Every challenge you overcome, every concept you master, and every connection you make is building the person you will become.

    \vspace{0.5em}

    \quotetext{"Education is the most powerful weapon which you can use to change the world." \\ --- Nelson Mandela}
\end{inspirational}

\subsection*{Your FUTA Legacy Awaits}

FUTA isn't just an institution --- it's an ecosystem of innovation, excellence, and transformation. Here, you'll:

\begin{itemize}[leftmargin=1.5em, itemsep=0.3em]
    \item Develop critical thinking and problem-solving skills
    \item Connect with future industry leaders
    \item Discover your passions and purpose
    \item Build resilience and character
    \item Create memories that will last a lifetime
\end{itemize}

\subsection*{What This Guide Offers}

\begin{itemize}[leftmargin=1.5em, itemsep=0.3em]
    \item Academic strategies for first-class performance
    \item Personal development insights
    \item Course-by-course battle plans
    \item Complete manual prices and buying guide
    \item A glossary of FUTA-specific terms
    \item Mental and spiritual wellness guidance
    \item Access to all course materials
    \item Complete registration guide
\end{itemize}

% ============================================================
% SECTION: THE MINDSET (NEW PAGE)
% ============================================================
\newpage
\section*{The Mindset}
\addcontentsline{toc}{section}{The Mindset}

\begin{center}
    {\color{primary}\Huge\bfseries Five Frameworks}\\[0.3em]
    {\color{textGray}\large To shape your FUTA journey --- from dream to reality}
\end{center}

\vspace{1em}

% ACTION FRAMEWORK
\begin{mindsetbox}
    \textbf{\color{primary}\Large ⚡ The ACTION Framework}\\[0.3em]
    {\color{textGray}\itshape How to turn dreams into reality}

    \vspace{1em}

    \textbf{Action} is a powerful word. For a dream to come true, \textbf{ACTION} is needed --- the same way \textbf{FORCE} causes motion.

    \vspace{1em}

    \begin{itemize}[leftmargin=1.5em, itemsep=0.5em]
        \item \textbf{\color{primary}A --- Aspire:} Have a clear vision and desire for your dream.
        \item \textbf{\color{primary}C --- Commit:} Dedicate yourself to making your dream a reality.
        \item \textbf{\color{primary}T --- Tenacity:} Persist and persevere, even in the face of obstacles.
        \item \textbf{\color{primary}I --- Initiative:} Take the first step and keep moving forward.
        \item \textbf{\color{primary}O --- Opportunity:} Recognize and seize opportunities that align with your dream.
        \item \textbf{\color{primary}N --- Navigate:} Adapt and adjust your approach as needed to stay on track.
    \end{itemize}

    \vspace{1em}

    By taking ACTION, you'll be well on your way to turning your dreams into reality. At first it is not easy, but remember that FORCE not only causes motion --- it is also the bridge from vision to mission.
\end{mindsetbox}

% FORCE FRAMEWORK
\begin{mindsetbox}
    \textbf{\color{primary}\Large 🔥 The FORCE Framework}\\[0.3em]
    {\color{textGray}\itshape Become unstoppable}

    \vspace{1em}

    \begin{itemize}[leftmargin=1.5em, itemsep=0.5em]
        \item \textbf{\color{primary}F --- Focus:} Concentrate your efforts and attention on your goals.
        \item \textbf{\color{primary}O --- Optimism:} Maintain a positive attitude and believe in yourself.
        \item \textbf{\color{primary}R --- Resilience:} Bounce back from setbacks and keep moving forward.
        \item \textbf{\color{primary}C --- Courage:} Summon the strength to take risks and face challenges head-on.
        \item \textbf{\color{primary}E --- Endurance:} Persevere through difficulties and stay committed to your objectives.
    \end{itemize}

    \vspace{1em}

    By applying FORCE to your endeavors, you'll be unstoppable. \textbf{Take ACTION and be FORCED.}
\end{mindsetbox}

% P5 PRINCIPLE
\begin{mindsetbox}
    \textbf{\color{primary}\Large 🎓 The P5 Principle}\\[0.3em]
    {\color{textGray}\itshape Proper Preparation Prevents Poor Performance --- plus the P that changes everything}

    \vspace{1em}

    Always remember \textbf{P5} --- not PS5.

    \vspace{0.5em}

    \begin{center}
        {\color{primary}\Large\bfseries Proper Preparation Prevents Poor Performance.}
    \end{center}

    \vspace{0.5em}

    In all you do, still apply \textbf{P}.

    \vspace{0.5em}

    \begin{successbox}
        You can't do without having \textbf{P} --- but \textbf{P} is the solution to \textbf{P}, and you'll be \textbf{P}.

        \vspace{0.5em}

        Make sure you \textbf{Pray}.

        \vspace{0.5em}

        You can't do without having \textbf{Problem}, but \textbf{Prayer} is a solution to \textbf{Problem}, and you'll be \textbf{Passed}.
    \end{successbox}

    \vspace{0.5em}

    {\color{primary}\itshape A word is enough for the wise.}
\end{mindsetbox}

% 6 vs 9
\begin{mindsetbox}
    \textbf{\color{primary}\Large 🔄 The 6 and 9 Metaphor}\\[0.3em]
    {\color{textGray}\itshape The difference between a dream and a reality}

    \vspace{1em}

    Academic excellence has been a dream for everyone in life. But do all dreams come true?

    \vspace{0.5em}

    This question has been asked many times. It's just like writing the figure \textbf{6}. You see it as 6 while others view it as 9. But 6 can turn into 9.

    \vspace{0.5em}

    \begin{inspirational}
        If you're reluctant, 6 will still remain 6 --- since you're not making an attempt to change it.

        \vspace{0.5em}

        But if you struggle to make an attempt, 6 will actually become 9.
    \end{inspirational}

    \vspace{0.5em}

    This means that for a dream to come true, \textbf{ACTION} is needed --- the same way \textbf{FORCE} causes motion.
\end{mindsetbox}

% ASK QUESTIONS
\begin{mindsetbox}
    \textbf{\color{primary}\Large ❓ Ask Questions}\\[0.3em]
    {\color{textGray}\itshape Nobody knows your mood}

    \vspace{1em}

    \begin{center}
        {\color{dangerRed}\Large\bfseries NOBODY KNOWS YOUR MOOD.}\\[0.3em]
        {\color{dangerRed}\Large\bfseries ASK QUESTIONS WHEN YOU FEEL CONFUSED.}
    \end{center}

    \vspace{1em}

    Don't be silent. Don't assume. Don't guess. Ask.

    \vspace{0.5em}

    Your lecturers, your seniors, your coursemates, FUTASUG, Student Affairs --- they're all here to help you. But they can only help if you open your mouth.
\end{mindsetbox}

% WORD TO FRESHERS
\begin{mindsetbox}
    \textbf{\color{primary}\Large 💙 A Word to Freshers}\\[0.3em]
    {\color{textGray}\itshape From the compiler}

    \vspace{1em}

    I want to congratulate you on your admission to the Federal University of Technology, Akure. This achievement is a testament to your hard work, dedication, and perseverance.

    \vspace{0.5em}

    I want to emphasize the importance of setting clear goals and priorities. Your time at FUTA is a precious opportunity to explore your passions, develop your skills, and build a strong foundation for your future. Be intentional about what you want to achieve, and make a plan to get there.

    \vspace{0.5em}

    FUTA offers a wide range of programs and courses that will challenge and enrich you. Take advantage of the resources available to you --- including your lecturers, classmates, and campus facilities. Don't be afraid to ask questions, seek help when needed, and explore new areas of interest.

    \vspace{0.5em}

    I also want to acknowledge that university life can be challenging. You may face obstacles, setbacks, and uncertainties along the way. But I urge you to remember that these experiences are an integral part of your growth and development.

    \vspace{0.5em}

    Don't be too hard on yourself when things don't go as planned. Instead, learn from your mistakes, and use them as opportunities for growth and improvement.

    \vspace{0.5em}

    Remember that you are not alone. Your fellow students, FUTASUG, lecturers, and Student Affairs are all here to help you succeed.
\end{mindsetbox}

% ============================================================
% SECTION 1: REGISTRATION (NEW PAGE)
% ============================================================
\newpage
\section{Complete Registration Guide}

Your registration journey begins when the portal opens. On this date, your FUTA journey officially begins. The registration process is your first major challenge, and this guide will help you navigate it successfully.

\begin{inspirational}
    \textbf{\color{primary}\Large 📅 Important Registration Timeline}

    \vspace{0.5em}

    \textbf{At the start of the session:} Undergraduate portal opens for online registrations and BIO-DATA FORM download

    \vspace{0.3em}

    \textbf{Shortly after:} Physical registration begins

    \vspace{0.5em}

    \quotetext{"Proper preparation prevents poor performance. Start your registration process early!"}
\end{inspirational}

\subsection{The BIO-DATA Form --- Your First Challenge}

The BIO-DATA form is the most tedious part of your registration. On the back page, you'll find six names with spaces for signatures:

\begin{enumerate}[leftmargin=1.5em, itemsep=0.3em]
    \item \textbf{Your School Officer} --- General officer who signs for all students
    \item \textbf{Your Registration Officer} --- Department-specific officer
    \item \textbf{Your Head of Department (HOD)} --- Head of your department
    \item \textbf{Your Dean of School} --- Head of your faculty
    \item \textbf{Dean of Students' Affairs Division}
    \item \textbf{The School Registrar}
\end{enumerate}

\textbf{Location Guide:} The first four signatures are from your faculty building, while the last two are outside your faculty.

\subsection{Step-by-Step Registration Process}

\begin{tcolorbox}[colback=lightGray,colframe=primary,title={\color{white}\textbf{Step 1: School Officer}},colbacktitle=primary,arc=4pt,breakable]
Your school officer is general and will sign the biodata of every student in the school. This is usually the quickest step. The officer handles all students across all departments in your school, so expect a queue --- but the process itself is fast. Have your form filled out correctly before you get to the front of the line to avoid being sent back. If you can, go early in the morning before the crowd builds up. Also, make sure your passport photograph is properly attached where required.
\end{tcolorbox}

\begin{tcolorbox}[colback=lightGray,colframe=primary,title={\color{white}\textbf{Step 2: Registration Officer}},colbacktitle=primary,arc=4pt,breakable]
Every department has its own registration officer. This person handles registration for your department specifically. They will verify that your courses are properly selected and that your bio-data matches departmental records. Since this is department-specific, the queue is usually shorter than the school officer's queue --- but the verification is more thorough. Bring all your documents here, and double-check every piece of information before they sign. Any errors here can cause problems later, so take your time.
\end{tcolorbox}

\begin{tcolorbox}[colback=lightGray,colframe=primary,title={\color{white}\textbf{Step 3: Head of Department (HOD)}},colbacktitle=primary,arc=4pt,breakable]
You should already know your HOD's name by now. If you don't, find out immediately! This is the head of your department --- the senior academic in charge. You'll need to see them personally for their signature. Dress well, greet respectfully, and be brief. Some HODs will sign on the spot; others may ask you to return at a specific time. Always confirm the HOD's office hours before going, and if there's a secretary, be polite to them too --- they can help you get in faster.
\end{tcolorbox}

\begin{tcolorbox}[colback=lightGray,colframe=primary,title={\color{white}\textbf{Step 4: Dean of School}},colbacktitle=primary,arc=4pt,breakable]
The Dean is the highest authority in your faculty. You might not see them physically during registration. Some faculties will collect all students' forms in one batch, present them to the Dean, and then distribute them back the next day. Other faculties require you to submit and return later to collect the signed form. Either way, don't panic --- just follow the instructions given at the Dean's office. The process is usually smoother than students expect because the Dean's office handles many submissions at once.
\end{tcolorbox}

\begin{tcolorbox}[colback=lightGray,colframe=primary,title={\color{white}\textbf{Step 5: Students' Affairs Division}},colbacktitle=primary,arc=4pt,breakable]
Once you have the Dean's signature, proceed to Students' Affairs. \textbf{Dress formally for this step!} Some staff in this office are strict about appearance and will send you back if you're not dressed properly. After conquering this part, you can wear whatever you like. Students' Affairs handles your welfare as a student --- they'll also want to see your documents one more time before signing. Be patient here; this is often the busiest office during registration.
\end{tcolorbox}

\begin{tcolorbox}[colback=lightGray,colframe=primary,title={\color{white}\textbf{Step 6: Final Submission}},colbacktitle=primary,arc=4pt,breakable]
Return to your faculty for final submission. You'll submit your file for the School Registrar. The Registrar is the highest administrative officer in the university, and their signature validates your registration. This is the last step --- once submitted, your registration is officially complete. Keep a photocopy of everything you submit for your own records, and make sure you collect any receipts or acknowledgment slips that are given.
\end{tcolorbox}

\subsection{Registration Rules \& Etiquette}

\begin{tcolorbox}[colback=lightGray,colframe=successGreen,title={\color{white}\textbf{👔 Dress Code}},colbacktitle=successGreen,arc=4pt,breakable]
\begin{itemize}[leftmargin=1em, itemsep=0.2em]
    \item Dress as neatly as possible for every registration step
    \item Have a good haircut or hairdo --- first impressions matter
    \item No need to overdress with expensive clothes; neatness is what counts
    \item When going to Students' Affairs, dress formally (shirt, trousers, or decent native wear for guys; appropriate formal or semi-formal for ladies)
    \item Avoid slippers, shorts, and casual flip-flops in administrative offices
\end{itemize}
\end{tcolorbox}

\begin{tcolorbox}[colback=lightGray,colframe=successGreen,title={\color{white}\textbf{🗣️ Communication}},colbacktitle=successGreen,arc=4pt,breakable]
\begin{itemize}[leftmargin=1em, itemsep=0.2em]
    \item Speak "most-politely" to all officials --- they deal with hundreds of students daily
    \item Listen to instructions carefully the first time to avoid repeats
    \item Ask questions respectfully when confused; don't argue
    \item Be patient and understanding --- queues can be long
    \item If you're sent back, don't take it personally; just fix the issue and return
    \item Greet first; "Good morning sir/ma" goes a long way
\end{itemize}
\end{tcolorbox}

\subsection{Financial Preparation}

\begin{moneybox}
    \textbf{\color{primary}\Large 💰 Recommended Budget: ₦20,000 (Cash or Bank Transfer)}

    \vspace{0.3em}

    \textit{Actual costs may vary by faculty.}

    \vspace{0.5em}

    \begin{itemize}[leftmargin=1.5em, itemsep=0.3em]
        \item \textbf{Faculty Dues} --- Payment to your faculty
        \item \textbf{Departmental Dues} --- Payment to your department
        \item \textbf{Engineering Students:} Blue Lab Coat
        \item \textbf{Faculty File} --- Official file for your faculty
        \item \textbf{Departmental File} --- Official file for your department
    \end{itemize}

    \vspace{0.5em}

    \textbf{📌 Important Reminders:}

    \begin{itemize}[leftmargin=1.5em, itemsep=0.3em]
        \item \textit{Carry both cash and digital payment options.} Some payments may only be accepted via bank transfer, while others may require cash.
        \item \textit{Confirm with your faculty for specific requirements.} Each department may have its own additional charges.
        \item \textit{Keep receipts for every payment.} Always collect and keep receipts --- they may be required for verification later.
    \end{itemize}
\end{moneybox}

\subsection{Document Preparation}

During physical registration, all your documents will be screened. If you're missing any required documents, you'll be sent back. Prepare all necessary documents in advance --- check the next section for the full document checklist.

\subsection{Pro Tips for Successful Registration}

\begin{inspirational}
    \begin{itemize}[leftmargin=1.5em, itemsep=0.3em]
        \item \textbf{Start early} --- Avoid the last-minute rush
        \item \textbf{Ask a senior} --- They know the process and can guide you
        \item \textbf{Be patient} --- Registration can take time, especially at busy offices
        \item \textbf{Double-check everything} --- Ensure all documents are complete before you leave home
        \item \textbf{Network with coursemates} --- Share information and help each other
    \end{itemize}
\end{inspirational}

% ============================================================
% SECTION 2: REQUIRED DOCUMENTS (NEW PAGE)
% ============================================================
\newpage
\section{Required Documents for Registration}

Prepare these documents in advance. Make at least 5 photocopies of each document.

\begin{inspirational}
    \textbf{\color{primary}\Large 📋 Document Preparation Guide}

    \vspace{0.5em}

    \textbf{Online Registration First:} Upload clear scans of all documents to the FUTA portal

    \vspace{0.3em}

    \textbf{Physical Registration Second:} Bring originals and photocopies for verification
\end{inspirational}

\subsection*{Complete Document Checklist}

\begin{tcolorbox}[colback=lightGray,colframe=primary,title={\color{white}\textbf{📄 FUTA Admission Letter}},colbacktitle=primary,arc=4pt,breakable]
Official admission letter from FUTA with your department and course details.

\textbf{Requirement:} \textcolor{dangerRed}{Physical Registration}

\begin{itemize}[leftmargin=1em, itemsep=0.2em]
    \item \textbf{How to get:} Download from student portal after paying acceptance fee
    \item \textbf{Acceptance fee:} ₦50,000 (previous session) --- confirm current amount from portal
    \item \textbf{Copies needed:} 5 copies
\end{itemize}
\end{tcolorbox}

\begin{tcolorbox}[colback=lightGray,colframe=primary,title={\color{white}\textbf{🎓 JAMB Admission Letter}},colbacktitle=primary,arc=4pt,breakable]
Official JAMB approval of your admission into FUTA.

\textbf{Requirement:} \textcolor{successGreen}{Online \& Physical}

\begin{itemize}[leftmargin=1em, itemsep=0.2em]
    \item \textbf{How to get:} Print from JAMB portal
    \item \textbf{Format:} 2-page document (Original \& Duplicate)
    \item \textbf{Use:} Duplicate for registration, Original for scholarships
\end{itemize}
\end{tcolorbox}

\begin{tcolorbox}[colback=lightGray,colframe=primary,title={\color{white}\textbf{🎓 Testimonial}},colbacktitle=primary,arc=4pt,breakable]
Secondary school leaving certificate verifying your conduct and attendance.

\textbf{Requirement:} \textcolor{dangerRed}{Physical Registration}

\begin{itemize}[leftmargin=1em, itemsep=0.2em]
    \item \textbf{How to get:} From your secondary school
    \item \textbf{Availability:} Depends on your former school
\end{itemize}
\end{tcolorbox}

\begin{tcolorbox}[colback=lightGray,colframe=primary,title={\color{white}\textbf{📊 UTME Result Slip}},colbacktitle=primary,arc=4pt,breakable]
Original JAMB score slip proving you met the cutoff mark.

\textbf{Requirement:} \textcolor{dangerRed}{Physical Registration}

\begin{itemize}[leftmargin=1em, itemsep=0.2em]
    \item \textbf{How to get:} Print from JAMB portal
    \item Primary evidence of meeting academic requirements
\end{itemize}
\end{tcolorbox}

\begin{tcolorbox}[colback=lightGray,colframe=primary,title={\color{white}\textbf{🏅 WAEC Certificate}},colbacktitle=primary,arc=4pt,breakable]
Original O'Level certificate (not result slip).

\textbf{Requirement:} \textcolor{successGreen}{Online \& Physical}

\begin{itemize}[leftmargin=1em, itemsep=0.2em]
    \item \textbf{How to get:} From your exam center/school
    \item \textbf{If unavailable:} Use Statement of Result
    \item \textbf{For two sittings:} Get both certificates; if not, use one certificate + statement of result for the second
\end{itemize}
\end{tcolorbox}

\begin{tcolorbox}[colback=lightGray,colframe=primary,title={\color{white}\textbf{✍️ Attestation Letter}},colbacktitle=primary,arc=4pt,breakable]
Character reference letter from authority figure.

\textbf{Requirement:} \textcolor{successGreen}{Online \& Physical}

\begin{itemize}[leftmargin=1em, itemsep=0.2em]
    \item \textbf{How to get:} From principal, religious leader, or community leader
    \item \textbf{Format:} Must be on letter-headed paper
\end{itemize}
\end{tcolorbox}

\begin{tcolorbox}[colback=lightGray,colframe=primary,title={\color{white}\textbf{⚖️ Court Affidavit}},colbacktitle=primary,arc=4pt,breakable]
Sworn declaration against secret cult membership.

\textbf{Requirement:} \textcolor{successGreen}{Online \& Physical}

\begin{itemize}[leftmargin=1em, itemsep=0.2em]
    \item \textbf{How to get:} Inside FUTA or any High Court
    \item \textbf{Process:} Fill information, get signature and red seal
\end{itemize}
\end{tcolorbox}

\begin{tcolorbox}[colback=lightGray,colframe=primary,title={\color{white}\textbf{🎂 Birth Certificate}},colbacktitle=primary,arc=4pt,breakable]
Official document declaring your date of birth.

\textbf{Requirement:} \textcolor{successGreen}{Online \& Physical}

\begin{itemize}[leftmargin=1em, itemsep=0.2em]
    \item \textbf{How to get:} From birth registry or declaration center
    \item Confirms minimum age requirement
\end{itemize}
\end{tcolorbox}

\begin{tcolorbox}[colback=lightGray,colframe=primary,title={\color{white}\textbf{📸 Passport Photographs}},colbacktitle=primary,arc=4pt,breakable]
Recent standardized passport photos.

\textbf{Requirement:} \textcolor{successGreen}{Online \& Physical}

\begin{itemize}[leftmargin=1em, itemsep=0.2em]
    \item \textbf{Specifications:} White or red background
    \item \textbf{Copies needed:} 20 recommended
    \item \textbf{Use:} Faculty/department records, library card, health center
\end{itemize}
\end{tcolorbox}

\begin{tcolorbox}[colback=lightGray,colframe=primary,title={\color{white}\textbf{📝 Bio-Data Form}},colbacktitle=primary,arc=4pt,breakable]
4-page form with personal details filled online.

\textbf{Requirement:} \textcolor{successGreen}{Online \& Physical}

\begin{itemize}[leftmargin=1em, itemsep=0.2em]
    \item \textbf{Process:} Fill online during registration
    \item \textbf{Important:} Page 4 must have signatures
    \item \textbf{Timeline:} Takes 2--5 days to complete
\end{itemize}
\end{tcolorbox}

\begin{tcolorbox}[colback=lightGray,colframe=primary,title={\color{white}\textbf{🆔 Local Government ID}},colbacktitle=primary,arc=4pt,breakable]
Certificate of identification from your LGA.

\textbf{Requirement:} \textcolor{successGreen}{Online \& Physical}

\begin{itemize}[leftmargin=1em, itemsep=0.2em]
    \item \textbf{Cruciality:} Highly important
    \item \textbf{Copies needed:} 5 copies sufficient
\end{itemize}
\end{tcolorbox}

\begin{inspirational}
    \textbf{\color{primary}\Large 💡 Document Assistance Available}

    \vspace{0.5em}

    I can help you obtain JAMB Admission Letter, UTME Result Slip, and WAEC Certificate at an affordable fee.

    \vspace{0.5em}

    \textbf{Abideen Yakub Opeyemi (O'yemi)} --- +234 705 037 6975
\end{inspirational}

% ============================================================
% SECTION 3: THE GREAT SHIFT (NEW PAGE)
% ============================================================
\newpage
\section{The Great Shift}

The first lesson of university life is self-leadership. You are now the CEO of your education and your future.

\begin{inspirational}
    \textbf{\color{primary}\Large 🚀 From Passenger to Pilot}

    \vspace{0.5em}

    In secondary school, you were following a set path. At FUTA, you're charting your own course. This shift from dependence to independence is your first and most important lesson as a university student.

    \vspace{0.5em}

    \quotetext{"The greatest day in your life is when you take total responsibility for yourself." \\ --- Denis Waitley}
\end{inspirational}

\subsection{Key Principles for Success}

\begin{tcolorbox}[colback=lightGray,colframe=primary,title={\color{white}\textbf{You Are in Charge}},colbacktitle=primary,arc=4pt,breakable]
Your lecturers are guides, not taskmasters. The drive to attend lectures, complete assignments, and delve deeper into topics must come from within. Nobody will chase you to class. Nobody will remind you about submissions. Nobody will check if you've studied. You're now responsible for your own learning --- and that's both a privilege and a responsibility. Embrace it fully. The students who thrive at FUTA are those who take ownership of their education from day one.
\end{tcolorbox}

\begin{tcolorbox}[colback=lightGray,colframe=primary,title={\color{white}\textbf{Master Your Time}},colbacktitle=primary,arc=4pt,breakable]
The freedom you now have is your greatest tool, or your biggest trap. You have more free time than ever before --- and nobody telling you what to do with it. Create a realistic schedule that balances academics, personal growth, and rest. Build routines. Learn to prioritize. Set aside time for lectures, study sessions, assignments, exercise, socializing, and sleep. A student with a plan will always outperform a student without one, even if the latter is naturally smarter.
\end{tcolorbox}

\begin{tcolorbox}[colback=lightGray,colframe=primary,title={\color{white}\textbf{Own Your Decisions}},colbacktitle=primary,arc=4pt,breakable]
Every choice --- from which lectures to attend, to how you spend your evenings, to who you surround yourself with --- shapes your FUTA experience and your future. Choose wisely, but also learn from mistakes. Nobody gets it perfect. The key is to learn fast, adjust, and keep moving forward. Don't be too hard on yourself when you slip up --- but don't ignore the lesson either. Every decision is a chance to build the person you're becoming.
\end{tcolorbox}

% ============================================================
% SECTION 4: COURSE SYNOPSIS (NEW PAGE)
% ============================================================
\newpage
\section{Course Synopsis (First Semester 100L)}

\begin{inspirational}
    \textbf{\color{primary}\Large 💙 Don't Be Scared --- You've Done Most of This Before!}

    \vspace{0.5em}

    Before you panic when you see the list of courses, breathe. Here's the truth: \textbf{most of your 100 level first semester courses are things you already learned in secondary school.} Biology, Chemistry, Physics, Mathematics --- they're just at a deeper level now. You're not starting from scratch; you're building on a foundation you already have.

    \vspace{0.5em}

    The only truly new things will be a few introductory courses that everyone takes (like Intro to Computing, Use of English, and Engineering Drawing). These are designed to \textit{introduce} you, not to fail you. Attend every class, take notes, do your assignments --- and you'll be fine.

    \vspace{0.5em}

    Also --- remember we use the word \textbf{"course"} here, not "subject". Get used to that. You'll take courses; you'll pass courses.
\end{inspirational}

\begin{warningbox}
    \textbf{\color{warningText}\Large ⚠️ IMPORTANT: These Are NOT All Your Courses}

    \vspace{0.5em}

    The courses listed in this guide are the \textbf{common/general courses} that most 100 level students take --- but they're \textbf{not the complete list} of your courses.

    \vspace{0.5em}

    Some courses are \textbf{specific to your department}. Some were \textbf{not yet published when we compiled this}. Some simply \textbf{didn't reach us in time}.

    \vspace{0.5em}

    \textbf{So don't assume this is everything.} You'll find the full list:

    \begin{itemize}[leftmargin=1.5em, itemsep=0.3em]
        \item On your \textbf{course form} during registration
        \item From your \textbf{class rep} or \textbf{departmental adviser}
        \item From \textbf{seniors in your department}
    \end{itemize}

    \vspace{0.5em}

    \textbf{Want to connect with your senior colleagues and departmental mates?} Message me directly --- I'll connect you to the right people in your department so you never miss a course.

    \vspace{0.5em}

    \textbf{Abideen Yakub Opeyemi (O'yemi)} --- +234 705 037 6975
\end{warningbox}

\subsection*{Important Note on Course Codes (CCMAS)}

In 2025/2026, FUTA switched from the \textbf{BMAS} curriculum to the \textbf{CCMAS} (Core Curriculum and Minimum Academic Standards) framework. This means some courses have new codes.

\textbf{Important --- Math codes are unchanged.} FUTA kept the \textbf{MTH} prefix. So:
\begin{itemize}[leftmargin=1.5em, itemsep=0.2em]
    \item \textbf{MTH 101} is still MTH 101 (not MTS)
    \item \textbf{MTH 102} is still MTH 102
    \item \textbf{MTH 104} is still MTH 104
\end{itemize}

\textbf{What changed:}
\begin{itemize}[leftmargin=1.5em, itemsep=0.2em]
    \item Some GNS courses became \textbf{GST} courses (e.g., GST 111 / GST 112)
    \item Some courses were merged, split, or restructured
    \item New courses were introduced in some departments
    \item Credit units may differ from what seniors took
\end{itemize}

\subsection*{Complete Course List with Synopsis}

% MTH 101
\begin{coursebox}
    \textbf{\color{primary}\Large MTH 101 --- Introductory Mathematics I}\\[0.3em]
    \textbf{Exam Type:} CBT/CBE \\
    \textbf{Manual:} Tutorial Manual --- ₦2,500

    \vspace{0.8em}

    \textbf{Full Synopsis:}

    Elementary set theory: subsets, union, intersections, complement, and Venn diagrams. Real number system and mathematical induction: integers, rational and irrational numbers, mathematical induction. Real sequences and series. Quadratic equation and Polynomial. Binomial theorem and expansion. Circular measure. Trigonometric functions.

    \vspace{0.8em}

    \textbf{\color{warningText}Warnings:}
    \begin{itemize}[leftmargin=1.5em, itemsep=0.2em]
        \item This is one of the courses you'll find interesting --- the lecturers always come to class.
        \item Do not miss any of the classes.
        \item Go to class early so you can sit at the front.
        \item You'll be doing tutorials organized by the school --- once a week. Do not dare to miss it.
    \end{itemize}
\end{coursebox}

% GST 111
\begin{coursebox}
    \textbf{\color{primary}\Large GST 111 --- Communication in English}\\[0.3em]
    \textit{(Formerly known as GNS 101 --- the old code is now GST 111 under CCMAS)}\\[0.3em]
    \textbf{Exam Type:} CBT/CBE \\
    \textbf{Manual:} ₦2,000

    \vspace{0.8em}

    \textbf{Full Synopsis:}

    GST 111 is English for Academic Purpose (EAP) course. The objective is to equip students with skills that are necessary for learning and studying effectively in a university and communicating in English as a Second Language.

    \textbf{Topics taught include:}
    \begin{enumerate}[leftmargin=1.5em, itemsep=0.2em]
        \item Time Management
        \item Study Skills
        \item Scientific Word Formation
        \item Parts of Speech
        \item Aspects of English Grammar
        \item Note-taking / Note-making
    \end{enumerate}

    \vspace{0.8em}

    \textbf{\color{warningText}Warnings:}
    \begin{itemize}[leftmargin=1.5em, itemsep=0.2em]
        \item Don't miss any of this class.
        \item There are class activities which might be tests and/or attendance which usually come unannounced.
        \item Apart from the physical class, you'll be given assessments to do online.
        \item All materials needed for this course will be downloaded individually. Ensure you download it yourself.
        \item For all the assessments you do online, there are marks for them. Ensure you do them all.
        \item Download your course materials yourself --- don't collect from your friends.
        \item You'll all be added to a Telegram group where lectures will be taking place.
        \item \textbf{PLEASE FOLLOW THIS COURSE LECTURER'S INSTRUCTIONS.} They are usually strict with their instructions.
    \end{itemize}
\end{coursebox}

% FGNS 103
\begin{coursebox}
    \textbf{\color{primary}\Large FGNS 103 --- Information Retrieval}\\[0.3em]
    \textit{(Formerly GNS 103 --- now FGNS 103 under CCMAS)}\\[0.3em]
    \textbf{Exam Type:} CBT/CBE \\
    \textbf{Manual:} ₦2,500

    \vspace{0.8em}

    \textbf{Full Synopsis:}

    Brief history of libraries; Library and education; University libraries and other types of libraries; Study skills (reference services); Types of library materials, using library resources including e-learning, e-materials, etc.; Understanding library catalogues (card, OPAC, etc.) and classification; Copyright and its implications; Database resources; Bibliographic citations and referencing. Development of modern ICT; Hardware technology; Software technology; Input devices; Storage devices; Output devices; Communication and internet services; Word processing skills (typing, etc.).

    \vspace{0.8em}

    \textbf{\color{warningText}Warnings:}
    \begin{itemize}[leftmargin=1.5em, itemsep=0.2em]
        \item You'll be given research assignments which must be done.
        \item You'll also be doing a library tour --- don't miss it. Lecturers don't joke with it.
        \item The library tour mark is 10--20 depending on your lecturer.
        \item There is sure a past question for this course.
    \end{itemize}
\end{coursebox}

% PHY 101
\begin{coursebox}
    \textbf{\color{primary}\Large PHY 101 --- General Physics I (Mechanics)}\\[0.3em]
    \textbf{Exam Type:} CBT/CBE \\
    \textbf{Manual:} ₦2,500

    \vspace{0.8em}

    \textbf{Full Synopsis:}

    \textbf{1. Quantities, Units and Dimensions:} Space and Time, frames of reference, Invariance of physical laws, relativity of simultaneity, relativity of time intervals, relativity of length, units and dimension; standards and units, unit consistency and conversions.

    \textbf{2. Kinematics and Vectors:} Vector addition, components of vectors, unit vectors, products of vectors. Displacement, time and average velocity, instantaneous velocity, average acceleration, motion with constant acceleration, freely falling bodies, position and velocity vectors, acceleration vector, projectile motion. Motion in a circle and relative velocity.

    \textbf{3. Fundamental Laws of Mechanics:} Forces and interactions, Newton's first law, Newton's second law, mass and weight, Newton's third law. Statics and dynamics: application of Newton's laws, dynamics of particles, frictional forces, dynamics of circular motion. Galilean invariance, universal gravitation, gravitational potential energy, elastic potential energy, conservative and non-conservative forces.

    \textbf{4. Work and Energy:} Kinetic energy and the work-energy theorem, power, momentum and impulse, conservation of momentum, collisions and momentum conservation, elastic collisions, centre of mass.

    \textbf{5. Rotational Dynamics and Angular Momentum:} Angular velocity and acceleration, energy in rotational motion, parallel axis theorem, torque, torque and rotation about a moving axis.

    \textbf{6. Simple Harmonic Motion and its Applications:} The simple pendulum, damped oscillations, forced oscillations and resonance.

    \vspace{0.8em}

    \textbf{\color{warningText}Warnings:}
    \begin{itemize}[leftmargin=1.5em, itemsep=0.2em]
        \item This is the same as your secondary school physics, but a little advanced. Read it very well.
        \item \textbf{CRITICAL: SHM (Simple Harmonic Motion) is NOT in the manual} --- but appeared in a test where every question was on SHM.
        \item \textbf{READ ALL THE TOPICS IN THIS COURSE SYNOPSIS} --- not only the ones in the manual bought.
        \item There are questions at the back of the manual --- study them.
    \end{itemize}
\end{coursebox}

% PHY 107
\begin{coursebox}
    \textbf{\color{primary}\Large PHY 107 --- Experimental Physics I}\\[0.3em]
    \textbf{Exam Type:} CBE (no CBT component) \\
    \textbf{Manual:} ₦2,500

    \vspace{0.8em}

    \textbf{Full Synopsis:}

    Exploratory practical course. Hands-on training in the use of laboratory equipment and report writing. Six major experiments:
    \begin{enumerate}[leftmargin=1.5em, itemsep=0.2em]
        \item Measurements
        \item Determination of acceleration due to gravity using spiral spring method
        \item Determination of acceleration due to gravity using compound pendulum method
        \item Determination of moment of inertia of a rigid body
        \item Determination of refractive index of glass using rectangular prism
        \item Determination of refractive index of glass using triangular prism
    \end{enumerate}

    \vspace{0.8em}

    \textbf{\color{warningText}Warnings:}
    \begin{itemize}[leftmargin=1.5em, itemsep=0.2em]
        \item You'll not be doing test and exam --- but there are online practicals on Saturday evenings.
        \item Online practicals are posted on the Physics department Facebook page.
        \item Buy the practical manual --- essential.
    \end{itemize}
\end{coursebox}

% CHM 101
\begin{coursebox}
    \textbf{\color{primary}\Large CHM 101 --- Basic General Chemistry I}\\[0.3em]
    \textbf{Exam Type:} CBT/CBE \\
    \textbf{Manual:} ₦3,000

    \vspace{0.8em}

    \textbf{Full Synopsis:}

    Atoms, atomic structures, atomic theory, aufbau method, Hund's rule, Pauli Exclusion principle, atomic spectra. Molecules and chemical reactions. Energetics. Chemical equation and stoichiometry. Atomic structure; Modern electronic theory of atoms. Radioactivity. Chemical kinetics, collision theory, Kinetic theory of gases. Solution, solubility and solubility product. Electrochemistry, electrode potential, half-cell equation.

    \vspace{0.8em}

    \textbf{\color{warningText}Warnings:}
    \begin{itemize}[leftmargin=1.5em, itemsep=0.2em]
        \item \textbf{Textbook: Fundamental Chemistry} --- this is where the questions are set. Get it.
        \item This is secondary school topics but a little advanced.
        \item Lecturers take attendance sometimes, especially when the class is not equipped.
        \item Read everything in the course synopsis.
    \end{itemize}
\end{coursebox}

% CHM 103 / CHM 107
\begin{coursebox}
    \textbf{\color{primary}\Large CHM 103 / CHM 107 --- Experimental Chemistry I}\\[0.3em]
    \textbf{Exam Type:} CBE (no CBT component) \\
    \textbf{Manual:} ₦2,500

    \vspace{0.8em}

    \textbf{Full Synopsis:}

    Introduction to basic laboratory procedure and apparatus in the chemistry laboratory. Recording of laboratory data. Calibration of basic laboratory equipment. Preparation and standardization of common reagents. Preparation of sulphide of metals and determination of their empirical formula. Determination of atomic weight of metals. Stoichiometry. Preparation of simple, double, and complex salts.

    \vspace{0.8em}

    \textbf{\color{warningText}Warnings:}
    \begin{itemize}[leftmargin=1.5em, itemsep=0.2em]
        \item This course is very easy to pass.
        \item Don't rely on the manual for your test and exam.
        \item Find a better textbook and read all the topics in this course synopsis.
    \end{itemize}
\end{coursebox}

% BIO 101
\begin{coursebox}
    \textbf{\color{primary}\Large BIO 101 --- General Biology I}\\[0.3em]
    \textbf{Exam Type:} CBT/CBE \\
    \textbf{Manual:} ₦2,500

    \vspace{0.8em}

    \textbf{Full Synopsis:}

    The scope of Biology and its place in human welfare including characteristics of life, concepts in biology, topical issues in biology and career opportunities. Diversity and classification of living things. Cell structure and organization; functions of cellular organelles; diversity, general reproduction, interrelationship of organisms, heredity and evolution; elements of ecology and types of habitat. Differences between plants and animals. Variation and life cycles of plants to include non-vascular plants like algae, fungi, bacteria, viruses, bryophytes and pteridophytes. Varieties and forms, life cycles and functions of flowering plants.

    \vspace{0.8em}

    \textbf{\color{warningText}Warnings:}
    \begin{itemize}[leftmargin=1.5em, itemsep=0.2em]
        \item This is the normal biology you did in your secondary school.
        \item You won't want to miss the class. It's usually fun.
        \item You'll be sent the materials by the lecturers in charge.
    \end{itemize}
\end{coursebox}

% BIO 103 / BIO 107
\begin{coursebox}
    \textbf{\color{primary}\Large BIO 103 / BIO 107 --- General Biology Practical I}\\[0.3em]
    \textbf{Exam Type:} CBE (no CBT component) \\
    \textbf{Manual:} ₦3,000

    \vspace{0.8em}

    \textbf{Full Synopsis:}

    Introduction to Biology practical use of the microscope, cell structure, staining starch in plant tissue, algae, fungi, bryophytes and liverwort, gymnosperms, angiosperms.

    \vspace{0.8em}

    \textbf{\color{warningText}Warnings:}
    \begin{itemize}[leftmargin=1.5em, itemsep=0.2em]
        \item Buy the manual.
        \item Attend practical class.
        \item Do your class works and assignments.
        \item You'll be doing test and exam on this course.
    \end{itemize}
\end{coursebox}

% CSC 101
\begin{coursebox}
    \textbf{\color{primary}\Large CSC 101 --- Introduction to Computer Science}\\[0.3em]
    \textbf{Exam Type:} CBT/CBE \\
    \textbf{Manual:} ₦3,000

    \vspace{0.8em}

    \textbf{Full Synopsis:}

    History of computers; functional components of a computer; characteristics of a computer system. Broad introduction to programming methodology and algorithm. Emphasis on problem-solving strategies and techniques for developing/documenting applications, including principles of structured programming, problem decomposition, programme organization, the use of procedural abstraction and basic debugging skills. Visual BASIC programming language serves as the vehicle to illustrate the many concepts.

    \vspace{0.8em}

    \textbf{\color{warningText}Warnings:}
    \begin{itemize}[leftmargin=1.5em, itemsep=0.2em]
        \item Not everyone will offer this course, but to whom it may concern: you'll be having a practical class and a theoretical class.
        \item You can miss the theoretical class, but don't joke with the practical class.
        \item Attendance will always be taken.
        \item You'll be buying the practical manual.
    \end{itemize}
\end{coursebox}

% FMEE 101
\begin{coursebox}
    \textbf{\color{primary}\Large FMEE 101 --- Engineering Drawing I}\\[0.3em]
    \textbf{Exam Type:} Written \\
    \textbf{Manual:} ₦3,000

    \vspace{0.8em}

    \textbf{Full Synopsis:}

    This course is designed primarily for all students admitted into the Federal University of Technology, Akure. It provides a comprehensive knowledge and insight into engineering drawing as a basic tool of engineering.

    \textbf{Topics to be covered include:}
    \begin{itemize}[leftmargin=1.5em, itemsep=0.2em]
        \item Instruments for engineering drawing and their uses
        \item Drawing Paper Sizes; Margins; and Title Blocks
        \item Lettering and types of line
        \item Geometrical construction: bisection of lines and angles and their applications
        \item Polygon, tangency, locus of simple mechanisms
        \item Pictorial drawing; Isometric, oblique and perspectives
        \item Orthographic projection
        \item Dimensioning and development of simple shapes
        \item Assembly drawing of simple components
        \item Conventional representation of common engineering features
        \item Freehand sketching
    \end{itemize}

    \vspace{0.8em}

    \textbf{\color{warningText}Warnings:}
    \begin{itemize}[leftmargin=1.5em, itemsep=0.2em]
        \item This course is the simplest course in FUTA --- but because some of us are lazy and don't like stress, we say it's very hard. You will enjoy it.
        \item Attend class. Do your class works and assignments.
        \item This course will want to chase you away --- but remember FORCE.
    \end{itemize}
\end{coursebox}

% GET 101 / MNE 201
\begin{coursebox}
    \textbf{\color{primary}\Large GET 101 / MNE 201 --- Engineers in Society}\\[0.3em]
    \textit{(The course code may be GET 101 or MNE 201 depending on your department --- check your course form.)}\\[0.3em]
    \textbf{Exam Type:} CBT/CBE \\
    \textbf{Manual:} \textbf{Free PDF} (no physical manual to buy)

    \vspace{0.8em}

    \textbf{Full Synopsis:}

    \textbf{Section A: Early Development of Technology} --- What is Technology? History of Technology. Prehistoric Period. Stone Age, Bronze Age, Iron Age. Early Contributors to Technology Development: The Egyptians Empire (3000 BC -- 500 BC), The Greek Empire (600 BC -- 400 AD), The Romans Civilization Period. History of Materials Discoveries and Uses: Discovery of Fire, Stone and Wood, Uses of Animals, Solid Minerals. Revolution of the Engineering Industries (First through Fifth Industrial Revolutions). History of Engineering in the World and in Nigeria (PWD, NSTDA, ITF, RMRDC, etc.).

    \textbf{Section B: Historical Development of Science, Engineering and Technology} --- Relationship between Science, Engineering, and Technology.

    \textbf{Section C: Development of Engineering Profession} --- Engineering Profession. Codes of Practice and Standards. Ethics in Engineering.

    \textbf{Section D: Engineering Regulation and Professional Bodies in Nigeria} --- COREN (Council for the Regulation of Engineering Practices in Nigeria). NSE (Nigerian Society of Engineers). COMEG (Council of Nigerian Mining Engineers and Geoscientists). Engineering Cadres (Engineer, Technologist, Technician, Craftsman). Engineering Registration.

    \textbf{Section E: Engineering Education / Training} --- Introduction. Rationale for Training Engineering Profession. Brief History. Challenges of Engineering Education in Nigeria.

    \textbf{Section F: Roles and Challenges of Engineering in National Building} --- Agricultural sector. Banking. Automobile. Educational sector. Marketing. Health sector. Community Development.

    \textbf{Section G: Health, Safety \& Environmental (HSE) Management in Engineering Practice} --- HSE management phases. Compliance with standards. Engineering Code of Practice. Safety engineering principles. Bad Engineering Practices. HSE Sustainability Management.

    \textbf{Section H: Engineering Risk Analysis} --- Risk management. Steps in Engineering Risk Analysis: Risk Identification, Risk Impact Assessment, Risk Prioritization, Risk Mitigation Planning. Risk Mitigation Strategies.

    \vspace{0.8em}

    \textbf{\color{warningText}Warnings:}
    \begin{itemize}[leftmargin=1.5em, itemsep=0.2em]
        \item This course is free --- no manual to buy. The course material will be given as a PDF.
        \item \textbf{German questions:} Questions where you supply accurate and rightly-spelt answers that meet the lecturer's satisfaction.
        \item You have to know every part of this course material.
        \item Questions come from any angle, but won't exceed the course material.
    \end{itemize}
\end{coursebox}

\subsection*{Departmental Courses}

In addition to the general courses above, some departments have specific courses that only their students take. These may include (but are not limited to):

\begin{itemize}[leftmargin=1.5em, itemsep=0.2em]
    \item \textbf{EEE 101} --- Electrical Engineering
    \item \textbf{MNE 101} --- Mining Engineering
    \item \textbf{MCS 101} --- Mathematical Sciences
    \item \textbf{BTH 101} --- Biotechnology
    \item \textbf{IFS 101} --- Information Systems
    \item \textbf{URP 101} --- Urban and Regional Planning
\end{itemize}

Check your course form during registration to see exactly which departmental courses apply to you. When in doubt, ask your class rep, your departmental adviser, or check the departmental notice boards.

% ============================================================
% SECTION 5: ACADEMIC PLAYBOOK (NEW PAGE)
% ============================================================
\newpage
\section{Your Academic Playbook}

Your number one goal is academic excellence. To achieve this, you need a strategy.

\begin{inspirational}
    \textbf{\color{primary}\Large 🎓 P5 Principle: Proper Preparation Prevents Poor Performance}

    \vspace{0.5em}

    Success in FUTA isn't about being the smartest --- it's about being the most prepared, consistent, and resilient. Your daily habits determine your semester results.

    \vspace{0.5em}

    \quotetext{"The will to win is nothing without the will to prepare." \\ --- Juma Ikangaa}
\end{inspirational}

\subsection*{Benefits of Attending Classes}

\begin{tcolorbox}[colback=lightGray,colframe=primary,title={\color{white}\textbf{Attendance Marks}},colbacktitle=primary,arc=4pt,breakable]
Some courses award marks for class attendance, which contributes to your final grade. This is essentially free marks --- you just need to show up. For some courses, attendance can be worth 5--10\% of your total grade. Missing classes casually means throwing those marks away. Don't do it.
\end{tcolorbox}

\begin{tcolorbox}[colback=lightGray,colframe=primary,title={\color{white}\textbf{Better Understanding}},colbacktitle=primary,arc=4pt,breakable]
Lecturers explain concepts in ways that textbooks can't. They give real-life examples, tell stories, and highlight what really matters for exams. You learn faster when you're in the room. Some lecturers even drop hints about what will come out in the test --- hints you would completely miss if you weren't there.
\end{tcolorbox}

\begin{tcolorbox}[colback=lightGray,colframe=primary,title={\color{white}\textbf{Ask Questions}},colbacktitle=primary,arc=4pt,breakable]
Get immediate clarification on confusing topics directly from the expert. Don't let a small misunderstanding compound over weeks until it becomes a big problem. Ask questions in class. If you're shy, approach the lecturer after class or during office hours. Most lecturers are happy to explain things further.
\end{tcolorbox}

\begin{tcolorbox}[colback=lightGray,colframe=primary,title={\color{white}\textbf{Test Structure}},colbacktitle=primary,arc=4pt,breakable]
Learn the exam format --- whether it's objectives, MCQs, or fill-in-the-gaps (German). Lecturers usually drop hints about this before the test. Knowing the format in advance helps you prepare appropriately. For CBT courses, you need speed and accuracy. For written exams, you need to show your workings and write legibly.
\end{tcolorbox}

\subsection*{Exam and Test Strategy}

\begin{tcolorbox}[colback=lightGray,colframe=primary,title={\color{white}\textbf{📝 Before Tests/Exams}},colbacktitle=primary,arc=4pt,breakable]
\begin{itemize}[leftmargin=1em, itemsep=0.3em]
    \item Start preparation early --- don't wait till the last minute
    \item Create a study schedule and stick to it
    \item Practice with past questions --- this is the single most effective strategy
    \item Get enough sleep the night before --- a tired brain doesn't think well
    \item Prepare all necessary materials in advance (pen, pencil, calculator, ID card)
    \item Review previous test questions from your friends --- patterns repeat
\end{itemize}
\end{tcolorbox}

\begin{tcolorbox}[colback=lightGray,colframe=primary,title={\color{white}\textbf{📅 On Test/Exam Day}},colbacktitle=primary,arc=4pt,breakable]
\begin{itemize}[leftmargin=1em, itemsep=0.3em]
    \item Arrive early to the venue --- don't rush
    \item Stay calm and focused --- take deep breaths if you feel anxious
    \item Read instructions carefully --- especially for CBT where you can accidentally skip
    \item Manage your time effectively --- don't spend too long on one question
    \item For CBT: Practice speed and accuracy --- the timer doesn't wait
    \item For written: Write legibly and show your workings --- partial marks are better than nothing
    \item Answer the easy questions first, then tackle the harder ones
\end{itemize}
\end{tcolorbox}

\begin{tcolorbox}[colback=lightGray,colframe=primary,title={\color{white}\textbf{📊 After Tests/Exams}},colbacktitle=primary,arc=4pt,breakable]
\begin{itemize}[leftmargin=1em, itemsep=0.3em]
    \item Immediately gather the questions you remember --- write them down or send them to your group chat
    \item These questions often reappear in the final exam, sometimes word-for-word
    \item Don't make the same mistake twice --- analyze where you struggled
    \item Review your performance and learn from mistakes
    \item Start preparing for the next assessment immediately --- momentum matters
    \item Share the questions with your coursemates --- helping others helps you remember too
\end{itemize}
\end{tcolorbox}

\begin{inspirational}
    \textbf{\color{primary}\Large 💡 Important Reminder}

    \vspace{0.5em}

    For both computer-based and written exams, immediately after your test, gather questions from your friends that you remember. These are likely to come out in the final exam. Don't miss this opportunity twice!
\end{inspirational}

% ============================================================
% SECTION 6: MANUALS & PRICES (NEW PAGE)
% ============================================================
\newpage
\section{Manuals \& Prices (2025/2026)}

\begin{warningbox}
    \textbf{\color{warningText}\Large 📌 Important Notice}

    \vspace{0.5em}

    Manual prices change from session to session. The prices below are for the \textbf{2025/2026} session. \textbf{Always confirm the current price from the school bookshop or your department before buying.}
\end{warningbox}

\vspace{1em}

\begin{center}
\begin{longtable}{p{2cm} p{5cm} p{2.5cm} p{6cm}}
    \toprule
    \textbf{\color{primary}Course} & \textbf{\color{primary}Manual} & \textbf{\color{primary}Price} & \textbf{\color{primary}Notes} \\
    \midrule
    \endhead

    MTH 101 & Tutorial Manual & \textbf{₦2,500} & Essential for weekly tutorials \\
    \midrule
    GST 111 & Course Manual & \textbf{₦2,000} & For the physical class activities \\
    \midrule
    FGNS 103 & Course Manual & \textbf{₦2,500} & Library tour and research assignments \\
    \midrule
    PHY 101 & Course Manual & \textbf{₦2,500} & Questions at the back --- study them \\
    \midrule
    PHY 103 & Course Manual & \textbf{₦2,500} & Assignments given from this manual \\
    \midrule
    PHY 107 & Practical Manual & \textbf{₦2,500} & For the 6 major experiments \\
    \midrule
    CHM 101 & Course Manual & \textbf{₦3,000} & Textbook: Fundamental Chemistry \\
    \midrule
    CHM 103 / 107 & Practical Manual & \textbf{₦2,500} & Lab procedure guide \\
    \midrule
    BIO 101 & Course Manual & \textbf{₦2,500} & Lecturer sends extra materials \\
    \midrule
    BIO 103 / 107 & Practical Manual & \textbf{₦3,000} & Microscope and lab work \\
    \midrule
    CSC 101 & Practical Manual & \textbf{₦3,000} & Attendance taken in practical class \\
    \midrule
    FMEE 101 & Practical Manual & \textbf{₦3,000} & Class exercises drawn in real time \\
    \midrule
    GET 101 / MNE 201 & Free PDF & \textbf{Free} & No physical manual --- PDF only \\
    \bottomrule
\end{longtable}
\end{center}

\begin{successbox}
    \begin{center}
        \textbf{\color{primary}\Large 💰 Estimated Total Cost}\\[0.5em]
        {\color{successGreen}\Huge\bfseries ₦32,500}\\[0.5em]
        \textit{Based on 12 paid manuals. 1 manual is free (GET 101 / MNE 201).}
    \end{center}
\end{successbox}

\subsection*{Money-Saving Tips}

\begin{inspirational}
    \begin{itemize}[leftmargin=1.5em, itemsep=0.3em]
        \item \textbf{Buy in groups:} Some courses have a shared manual that multiple students can read from --- split the cost with coursemates.
        \item \textbf{Buy from seniors:} Often cheaper, sometimes pre-annotated with study notes.
        \item \textbf{Don't skip free manuals:} GET 101 / MNE 201 has a free PDF --- don't pay for a printed version unless you really want one.
        \item \textbf{Budget early:} Save up the total cost before resuming, so you're not caught off guard.
    \end{itemize}
\end{inspirational}

% ============================================================
% SECTION 7: BEYOND THE BOOKS (NEW PAGE)
% ============================================================
\newpage
\section{Beyond the Books}

FUTA offers more than a certificate. It offers an experience that will shape your character and future.

\begin{inspirational}
    \textbf{\color{primary}\Large 🌟 Your Education Extends Beyond the Classroom}

    \vspace{0.5em}

    The most valuable lessons often happen outside lecture halls --- in student organizations, late-night discussions with roommates, and challenges you overcome together. Your degree matters, but the person you become while earning it matters even more.

    \vspace{0.5em}

    \quotetext{"The beautiful thing about learning is that no one can take it away from you." \\ --- B.B. King}
\end{inspirational}

\subsection{Building Your Complete FUTA Experience}

\begin{tcolorbox}[colback=lightGray,colframe=primary,title={\color{white}\textbf{Network Intentionally}},colbacktitle=primary,arc=4pt,breakable]
The friends you make here will be your professional network for life. Connect with people who challenge and inspire you. Join study groups. Attend departmental events. Get to know your coursemates beyond the surface level. Some of the most important connections you'll ever make happen in university --- don't sleep on them.
\end{tcolorbox}

\begin{tcolorbox}[colback=lightGray,colframe=primary,title={\color{white}\textbf{Get Involved}},colbacktitle=primary,arc=4pt,breakable]
Join a club, society, or sports team. It builds character, leadership, and lifelong friendships. Don't be a ghost --- be present. Whether it's a departmental association, a religious group, or a hobby club, being part of something bigger than yourself enriches your experience. You'll learn leadership, teamwork, and how to manage multiple responsibilities --- all skills that will serve you long after graduation.
\end{tcolorbox}

\begin{tcolorbox}[colback=lightGray,colframe=primary,title={\color{white}\textbf{Prioritize Your Well-being}},colbacktitle=primary,arc=4pt,breakable]
Your mental and physical health are critical. Eat well, exercise, sleep, and seek help when needed. A healthy mind learns more efficiently than an exhausted one. If you're struggling with stress, anxiety, or depression, don't be silent --- talk to someone you trust, visit the health center, or reach out to a counsellor. Taking care of yourself isn't a luxury; it's a necessity.
\end{tcolorbox}

\begin{tcolorbox}[colback=lightGray,colframe=primary,title={\color{white}\textbf{Develop Soft Skills}},colbacktitle=primary,arc=4pt,breakable]
Communication, teamwork, and time management are as important as your technical knowledge. Employers look for people who can work with others, express ideas clearly, and handle pressure gracefully. Use your time at FUTA to develop these skills --- participate in group projects, take on leadership roles, and practice public speaking. These soft skills will often matter more than your CGPA when you enter the workforce.
\end{tcolorbox}

% ============================================================
% SECTION 8: NAVIGATING THE STORMS (NEW PAGE)
% ============================================================
\newpage
\section{Navigating the Storms}

There will be difficult days. Challenging courses, overwhelming workloads, and personal struggles are part of the journey.

\begin{inspirational}
    \textbf{\color{primary}\Large ⚡ Storms Make Better Sailors}

    \vspace{0.5em}

    Every successful FUTA graduate has faced challenges. What separates them is not the absence of storms, but their ability to navigate through them. You will face difficult moments --- a test you failed, a course you can't understand, a season of loneliness or doubt. These are not signs that you don't belong. They are signs that you're growing.

    \vspace{0.5em}

    \quotetext{"The oak fought the wind and was broken, the willow bent when it must and survived." \\ --- Robert Jordan}
\end{inspirational}

\subsection{Your Support System}

\begin{tcolorbox}[colback=lightGray,colframe=primary,title={\color{white}\textbf{FUTA Student Union Govt (FUTASUG)}},colbacktitle=primary,arc=4pt,breakable]
Your representatives are here to help with academic and welfare issues. Don't hesitate to reach out to them. From issues with lecturers to accommodation problems to general student welfare --- FUTASUG exists to support students. Get to know your departmental and faculty representatives; they can be powerful allies when you face challenges.
\end{tcolorbox}

\begin{tcolorbox}[colback=lightGray,colframe=primary,title={\color{white}\textbf{Departmental Associations \& Seniors}},colbacktitle=primary,arc=4pt,breakable]
Connect with seniors who've walked your path. They have the answers, the past questions, and the experience. Don't be shy --- most seniors are happy to help if you approach them respectfully. Your departmental association is also a great source of support, materials, and mentorship. When you're confused about a course, a lecturer, or the system in general, seniors are often the fastest source of clarity.
\end{tcolorbox}

\begin{tcolorbox}[colback=lightGray,colframe=primary,title={\color{white}\textbf{Lecturers \& Course Advisors}},colbacktitle=primary,arc=4pt,breakable]
Most lecturers are willing to help students who show genuine interest. Attend office hours. Ask intelligent questions. Show effort. Course advisors especially are there to guide you through academic decisions. If you're struggling with a course, the lecturer's office is the right place to be. Don't wait until the exam is near to seek help --- show up early and consistently.
\end{tcolorbox}

\begin{tcolorbox}[colback=lightGray,colframe=primary,title={\color{white}\textbf{Religious Groups}},colbacktitle=primary,arc=4pt,breakable]
The heads in your religious bodies are also there to help you. Balance your spiritual life with your academic life --- both matter. Whether you're Christian, Muslim, or of another faith, being part of a religious community on campus provides moral support, mentorship, and a sense of belonging. Campus fellowships often have senior members who can guide you through academic challenges too.
\end{tcolorbox}

\subsection{Turning Setbacks into Comebacks}

\textbf{Failure is Feedback:} A low score is not a label; it's data. Analyze it, learn from it, and come back stronger. Every top student has failed at something --- what matters is what they learned from it. When you don't get the grade you wanted, don't hide from it. Look at what went wrong. Adjust your approach. Try again. The students who succeed in the long run are not the ones who never fail --- they're the ones who refuse to stay down.

% ============================================================
% SECTION 9: THE ANCHOR (NEW PAGE)
% ============================================================
\newpage
\section{The Anchor: The Power of Prayer}

In the midst of the pressure, never underestimate the power of faith and prayer.

\begin{inspirational}
    \textbf{\color{primary}\Large 🙏 Your Spiritual "P" --- The Ultimate Game Changer}

    \vspace{0.5em}

    While you're studying hard and preparing thoroughly, remember that some battles are won on your knees. Your spiritual foundation will sustain you when everything else feels unstable. Faith isn't an escape from reality --- it's the anchor that holds you steady while you navigate reality.

    \vspace{0.5em}

    \quotetext{"You can't do without having Problems, but Prayer is a solution to Problem, and you'll be Passed."}
\end{inspirational}

\subsection{Integrating Faith into Your FUTA Journey}

\begin{tcolorbox}[colback=lightGray,colframe=primary,title={\color{white}\textbf{Start Your Day Right}},colbacktitle=primary,arc=4pt,breakable]
Begin each day with prayer and meditation to set the right tone. A calm, focused morning leads to a productive, peaceful day. Even if it's just 5 minutes --- commit to starting your day with intentional spiritual focus. You'll notice the difference in how you handle challenges when your foundation is solid.
\end{tcolorbox}

\begin{tcolorbox}[colback=lightGray,colframe=primary,title={\color{white}\textbf{Faith During Challenges}},colbacktitle=primary,arc=4pt,breakable]
When facing difficult courses or personal struggles, turn to prayer for strength and guidance. You are not alone. The same God who brought you this far will see you through. Prayer doesn't always change your circumstances, but it changes you --- it gives you peace, clarity, and the strength to keep going when you feel like giving up.
\end{tcolorbox}

\begin{tcolorbox}[colback=lightGray,colframe=primary,title={\color{white}\textbf{Gratitude \& Balance}},colbacktitle=primary,arc=4pt,breakable]
Regularly give thanks for the opportunity to study at FUTA. Remember that "faith without works is dead" --- pray like it depends on God, work like it depends on you. Both matter equally. Faith gives you the strength and perspective; work gives you the results. Don't neglect either one. Find the balance, and you'll thrive academically and spiritually.
\end{tcolorbox}

\subsection{Your Final Charge, Class of 2029}

You are the architects of the future. The world is waiting for the solutions you will design, the technologies you will build, and the leadership you will provide. It all starts here, at FUTA.

Take relentless \textbf{ACTION}. Be an unstoppable \textbf{FORCE}. Apply the \textbf{P5} principle daily. And anchor it all in Prayer.

\begin{center}
    {\Large\bfseries\color{primary} Your time is now. Make it legendary.}
\end{center}

% ============================================================
% SECTION 10: COURSE MATERIALS (NEW PAGE)
% ============================================================
\newpage
\section{Course Materials}

All course materials are available on Google Drive --- past questions, lecture slides, practical manuals, and textbooks organized by course.

\begin{tcolorbox}[colback=lightGray, colframe=primary, arc=4pt]
    \textbf{\color{primary}\Large 📁 Access Main Materials Folder}

    \vspace{0.5em}

    \url{https://drive.google.com/drive/folders/11-S3HppUkfATQhO_f27idoZsY6xwNKAQ}
\end{tcolorbox}

\subsection*{Available Course Materials}

Each course has its own dedicated folder with organized materials. Use the main folder link above to browse all courses.

% ============================================================
% SECTION 11: DEPARTMENTAL COURSES (NEW PAGE)
% ============================================================
\newpage
\section{Departmental Courses Information}

The courses listed in this guide are the \textbf{general courses} that almost all FUTA students take in their first year. Depending on your department, you may also have specific departmental courses.

\begin{inspirational}
    \textbf{\color{primary}\Large 🎯 Important Notice About Your Courses}

    \vspace{0.5em}

    \textbf{Some departments start their specialized courses from 100 level!} While the courses in this guide are essential for all students, you need to check your specific departmental requirements.
\end{inspirational}

\subsection*{How to Know Your Exact Courses}

\begin{itemize}[leftmargin=1.5em, itemsep=0.3em]
    \item \textbf{Check your course form during course registration} --- Carefully review all courses assigned to you during online registration
    \item \textbf{Consult 200 Level Students:} Reach out to seniors in your department --- they've been through this before
    \item \textbf{Departmental Notice Boards:} Check for course information posted by your department
    \item \textbf{Ask Your Course Adviser:} Each department has course advisers for fresh students
\end{itemize}

\begin{inspirational}
    \textbf{\color{primary}\Large 💡 Take Your Departmental Courses Seriously!}

    \vspace{0.5em}

    Your departmental courses form the foundation of your specialized knowledge. Don't make the mistake of focusing only on general courses while neglecting your core departmental subjects. These courses:

    \begin{itemize}[leftmargin=1.5em, itemsep=0.3em]
        \item Build your technical foundation
        \item Introduce you to your field of study
        \item Prepare you for advanced courses
        \item Often have higher credit units
    \end{itemize}

    \quotetext{"The expert in anything was once a beginner. Start strong with your departmental courses!"}
\end{inspirational}

\subsection*{Need Help?}

If you're unsure about your departmental courses or who to ask:

\begin{itemize}[leftmargin=1.5em, itemsep=0.3em]
    \item \textbf{Check your course form during registration}
    \item \textbf{Reach out to me directly} --- I can help connect you with the right people
    \item \textbf{Join your departmental WhatsApp group} --- this is where important information is shared
    \item \textbf{Visit your department office} --- the administrative staff can guide you
    \item \textbf{Attend departmental orientation} --- don't miss this important event
\end{itemize}

\begin{center}
    \textbf{\color{primary} Remember: Your first year sets the tone for your entire FUTA journey. Start strong, ask questions, and build good academic habits from day one!}
\end{center}

% ============================================================
% SECTION 12: GLOSSARY (NEW PAGE)
% ============================================================
\newpage
\section{Glossary of FUTA Terms}

FUTA has its own vocabulary. This glossary explains the terms you'll hear every day --- so you're never confused.

\subsection*{Exams}

\begin{tcolorbox}[colback=glossaryBg,colframe=accent,arc=4pt,breakable]
    \textbf{\color{primary}\large German Questions}\\[0.3em]
    A FUTA-specific term for exam questions where students must supply accurate and rightly-spelt answers that meet the lecturer's satisfaction. Unlike multiple-choice or matching questions, you write out the full answer.

    \vspace{0.5em}

    \textit{Example:} A question that says "Define photosynthesis" and expects a full written definition with correct spelling.
\end{tcolorbox}

\begin{tcolorbox}[colback=glossaryBg,colframe=accent,arc=4pt,breakable]
    \textbf{\color{primary}\large CBT}\\[0.3em]
    Computer-Based Test. A test taken on a computer, usually with multiple-choice questions. Speed and accuracy matter because the timer runs continuously.

    \vspace{0.5em}

    \textit{Example:} Most 100L courses have CBT tests --- practice with timed sessions to build speed.
\end{tcolorbox}

\begin{tcolorbox}[colback=glossaryBg,colframe=accent,arc=4pt,breakable]
    \textbf{\color{primary}\large CBE}\\[0.3em]
    Computer-Based Examination. Similar to CBT but typically for end-of-semester exams. Often practical-based.

    \vspace{0.5em}

    \textit{Example:} PHY 107 and BIO 103 use CBE --- practical content is tested via computer.
\end{tcolorbox}

\subsection*{Academics}

\begin{tcolorbox}[colback=glossaryBg,colframe=accent,arc=4pt,breakable]
    \textbf{\color{primary}\large CCMAS}\\[0.3em]
    Core Curriculum and Minimum Academic Standards. The national curriculum FUTA adopted in 2025/2026, replacing BMAS. It restructured some course codes and content across Nigerian universities.

    \vspace{0.5em}

    \textit{Example:} Some courses changed codes --- for example, GNS 101 became GST 111. MTH codes stayed the same.
\end{tcolorbox}

\begin{tcolorbox}[colback=glossaryBg,colframe=accent,arc=4pt,breakable]
    \textbf{\color{primary}\large CGPA}\\[0.3em]
    Cumulative Grade Point Average. The running average of all your grades from 100 level to your current level. This is what determines your degree classification.

    \vspace{0.5em}

    \textit{Example:} A CGPA of 4.50--5.00 puts you in First Class range.
\end{tcolorbox}

\begin{tcolorbox}[colback=glossaryBg,colframe=accent,arc=4pt,breakable]
    \textbf{\color{primary}\large GPA}\\[0.3em]
    Grade Point Average. The average of your grades for a single semester. Your CGPA is the running average of all your GPAs.

    \vspace{0.5em}

    \textit{Example:} Strong GPAs every semester compound into a strong CGPA.
\end{tcolorbox}

\begin{tcolorbox}[colback=glossaryBg,colframe=accent,arc=4pt,breakable]
    \textbf{\color{primary}\large SIWES}\\[0.3em]
    Students' Industrial Work Experience Scheme. A mandatory internship program for engineering and technology students, usually done in 300 level or 400 level.

    \vspace{0.5em}

    \textit{Example:} You'll spend several months in an industry to gain practical experience.
\end{tcolorbox}

\subsection*{Campus Life}

\begin{tcolorbox}[colback=glossaryBg,colframe=accent,arc=4pt,breakable]
    \textbf{\color{primary}\large HOD}\\[0.3em]
    Head of Department. The senior academic in charge of your department. You'll need their signature during registration.

    \vspace{0.5em}

    \textit{Example:} Your HOD signs the back page of your BIO-DATA form.
\end{tcolorbox}

\begin{tcolorbox}[colback=glossaryBg,colframe=accent,arc=4pt,breakable]
    \textbf{\color{primary}\large FUTASUG}\\[0.3em]
    Federal University of Technology Akure Student Union Government. The elected student body that represents student interests on campus.

    \vspace{0.5em}

    \textit{Example:} Reach out to FUTASUG for academic and welfare issues.
\end{tcolorbox}

\subsection*{Mindset}

\begin{tcolorbox}[colback=glossaryBg,colframe=accent,arc=4pt,breakable]
    \textbf{\color{primary}\large P5}\\[0.3em]
    Proper Preparation Prevents Poor Performance. The study principle that success comes from consistency and preparation, not last-minute effort.

    \vspace{0.5em}

    \textit{Example:} Start reading from week one --- P5 in action.
\end{tcolorbox}

\begin{tcolorbox}[colback=glossaryBg,colframe=accent,arc=4pt,breakable]
    \textbf{\color{primary}\large FORCE}\\[0.3em]
    Focus, Optimism, Resilience, Courage, Endurance. A framework for becoming unstoppable in your academic journey.

    \vspace{0.5em}

    \textit{Example:} When MEE 101 tries to chase you away --- remember FORCE.
\end{tcolorbox}

\begin{tcolorbox}[colback=glossaryBg,colframe=accent,arc=4pt,breakable]
    \textbf{\color{primary}\large ACTION}\\[0.3em]
    Aspire, Commit, Tenacity, Initiative, Opportunity, Navigate. A framework for turning dreams into reality.

    \vspace{0.5em}

    \textit{Example:} Take ACTION daily --- one small step toward your goals.
\end{tcolorbox}

% ============================================================
% SECTION 13: GENIUS POINT ACADEMY (NEW PAGE)
% ============================================================
\newpage
\section{Genius Point Academy}

\begin{gpabox}
    {\color{accent}\huge\bfseries 🎓 Build a Solid First-Year Foundation With Us}

    \vspace{0.7em}

    {\color{white}Starting university can feel overwhelming --- new courses, new environment, new expectations. At \textbf{Genius Point Academy}, we make the transition smooth and set you up for long-term academic success.}

    \vspace{0.7em}

    {\color{accent}\large\bfseries Our Goal For You:}

    \vspace{0.3em}

    {\color{white}
    \begin{itemize}[leftmargin=1.5em, itemsep=0.3em]
        \item \textbf{Excel in your 100L courses} --- understand concepts deeply, not just pass exams
        \item \textbf{Build a strong GPA from day one} --- because a strong CGPA starts in year one
        \item \textbf{Master exam technique} --- past questions, timing, and CBT practice
        \item \textbf{Gain confidence} --- walk into every test knowing you've prepared well
        \item \textbf{Get mentorship} --- learn from seniors who've been there and excelled
    \end{itemize}}

    \vspace{0.5em}

    {\color{accent}\large\bfseries What We Offer:}

    \vspace{0.3em}

    {\color{white}
    \begin{itemize}[leftmargin=1.5em, itemsep=0.3em]
        \item Tutoring for all 100L courses (MTH, PHY, CHM, BIO, GNS, etc.)
        \item Step-by-step breakdown of difficult topics
        \item Past question review and exam-focused practice
        \item Small-group sessions or one-on-one (your choice)
        \item Ongoing mentorship beyond the classroom
    \end{itemize}}

    \vspace{0.7em}

    {\color{white}We're a team of dedicated FUTA students and recent graduates who have walked the path you're on. \textbf{O'yemi} --- the author of this guide --- is also a member of the Genius Point Academy team.}

    \vspace{0.5em}

    {\color{accent}\large\bfseries Ready to Start?}

    \vspace{0.3em}

    {\color{white}Reach out directly:}

    \vspace{0.3em}

    {\color{accent}\textbf{\Large 📱 +234 705 037 6975}}

    \vspace{0.3em}

    {\color{white}\small WhatsApp: Mention "FUTA Survival Guide" when you message us.}
\end{gpabox}

% ============================================================
% SECTION 14: CONTACT (NEW PAGE)
% ============================================================
\newpage
\section{Contact \& Support}

\begin{inspirational}
    \textbf{\color{primary}\Large 💪 Believe in Yourself!}

    \vspace{0.5em}

    No course is too difficult if you focus, stay attentive, and stay motivated. You have the ability to succeed; just put in consistent effort, stay positive, and keep your eyes on your goals.

    \vspace{0.5em}

    \quotetext{"You can do it! The journey of a thousand miles begins with a single step."}
\end{inspirational}

\subsection*{Your Senior Fellow FUTArian}

\begin{center}
    {\Large\bfseries\color{primary} Abideen Yakub Opeyemi (Ó'YẸMÍ)}

    \vspace{0.5em}

    \faGraduationCap\ Department of Electrical and Electronics Engineering

    \vspace{0.3em}

    \faPhone\ \textbf{+234 705 037 6975}

    \vspace{0.3em}

    \faEnvelope\ \textbf{abideenyakub08@gmail.com}
\end{center}

\subsection*{How I Can Help You}

\begin{itemize}[leftmargin=1.5em, itemsep=0.3em]
    \item \textbf{Course clarification:} Confused about a topic?
    \item \textbf{Material requests:} Need specific past questions or textbooks?
    \item \textbf{General guidance:} Questions about registration, accommodation, or campus life?
    \item \textbf{Prayer support:} Sometimes we just need someone to pray with us.
\end{itemize}

\subsection*{I'm Also Part of Genius Point Academy}

Alongside this guide, I'm a tutor and team member at \textbf{Genius Point Academy} --- where we help 100L students build a solid foundation and crush their first year.

Want to boost your GPA, master difficult courses, and get mentored? Reach out at \textbf{+234 705 037 6975} and mention "FUTA Survival Guide."

% ============================================================
% FINAL PAGE
% ============================================================
\newpage

\begin{center}
    \vspace*{3cm}

    {\color{primary}\Huge\bfseries Thank You for Reading}

    \vspace{1cm}

    {\color{accent}\rule{8cm}{2pt}}

    \vspace{1cm}

    {\Large\itshape You are a FUTArian.\\
    Your journey starts now.\\
    Make it legendary.}

    \vspace{2cm}

    {\color{textGray}\small This guide is an independent student-created resource.\\
    It is not officially affiliated with FUTA.\\
    For guidance purposes only.}

    \vfill

    {\color{primary}\copyright\ 2026 FUTA Survival Guide\\
    Compiled by Abideen Yakub Opeyemi (Ó'YẸMÍ)}
\end{center}

\end{document}
